\section{Compression screening and router details}
\label{app:compression}
\suppressfloats[t]

Before the pipeline study, most of the effort went into the model.
\Cref{tab:compression} summarizes the screening, all of which used a
1,024-image development split of COCO train2017, disjoint from val2017,
with official COCOeval. Each method received the same short recovery
budget: 1,024 training images, five epochs and one seed.

\begin{table}[b]
\centering
\caption{Compression screening for YOLO26n (development split of 1,024
train2017 images, one seed, five recovery epochs, TensorRT 10.3). The two
blocks come from different measurement contexts and are not directly
comparable. Block (a): median host-measured copy+execute+copy per image
in a sequential application with \texttt{opt0} builds, which includes
the down-clocking of \cref{sec:diagnosis}. Block (b): back-to-back
\texttt{trtexec} GPU compute with \texttt{opt3} builds (driver keeps the
GPU near 1,020\,MHz; AP$_S$ not recorded in that harness).}
\label{tab:compression}
\small
\setlength{\tabcolsep}{3pt}
\begin{tabular}{lrrr}
\toprule
Variant & AP & AP$_S$ & ms \\
\midrule
\multicolumn{4}{l}{\emph{(a) sequential application, \texttt{opt0} build}} \\
FP16, fine-tuned (reference) & 0.4691 & 0.2164 & 21.8 \\
Mixed QAT (50/99 conv INT8) & 0.4532 & 0.2038 & 24.4 \\
L1 channel pruning (-5.6\% params) & 0.3741 & 0.1426 & 21.8 \\
Object-risk pruning (-5.6\% params) & 0.4007 & 0.1613 & 21.8 \\
L1 pruning + QAT & 0.3612 & 0.1372 & 24.3 \\
Object-risk pruning + QAT & 0.3924 & 0.1594 & 24.5 \\
\midrule
\multicolumn{4}{l}{\emph{(b) back-to-back \texttt{trtexec}, \texttt{opt3} build}} \\
FP16 & 0.4686 & -- & 4.07 \\
Mixed QAT & 0.4593 & -- & 5.00 \\
\bottomrule
\end{tabular}
\end{table}

Channel pruning removed 25\% of the output channels of four backbone
convolutions, or 5.6\% of the parameters. The proposed object-risk
criterion scores channels by gradient-weighted foreground activations
and a fake-INT8 perturbation proxy. At the same budget it kept 2.7 AP
points more than L1 magnitude pruning, but both pruned models still lost
6.8--9.5 AP points. Neither was faster after compilation, because
TensorRT does not spend its time in the pruned early convolutions.
Ultralytics QAT produced a mixed-precision engine (about half of the
convolutions in INT8, the rest FP16/FP32 behind reformatting layers). It
was slower than FP16 under both builder levels. Distillation from the
fine-tuned teacher reduced AP in both pruning settings (L1: 0.369 to
0.338; object-risk: 0.395 to 0.373 on the development split). Pruning
the parts that per-layer profiles flagged as expensive (P3/P4
classification-head channels and the attention block's feed-forward
width) gave no stable change in engine time beyond run-to-run noise.
That engine time is about 6.5\,ms back-to-back with \texttt{opt0}
builds. The table's 21.8\,ms column is the same engine inside a
sequential application, and the difference between the two is the
governor effect of \cref{sec:diagnosis}. These results come from
the short budget above, with \texttt{opt0} builds on one device. Longer
training would likely let the pruned models recover more AP, and a full
\texttt{opt3} build sweep might redistribute the engine time. The
results therefore do not show that compression
cannot work. They show that, for a 2.6\,M-parameter NMS-free detector
on this device and budget, the achievable engine-time savings are a
small fraction of the 3$\times$ governor effect of \cref{sec:diagnosis}.

\subsection{The per-image resolution router}
\label{app:router}

A logistic router was fitted on cheap pre-inference image statistics.
It chooses a 512 or 640 engine per image, with its threshold fixed to
send 20\% of images to 640. On its 512-image held-out half, it slightly
beat a fixed 576 engine in both AP and median time. It was then frozen
and evaluated once on an independent, hash-verified 1,024-image
train2017 split that played no part in fitting or threshold selection
(\cref{sec:setup}). Against fixed 576 it lost 0.0089 AP (0.4386 vs.\
0.4476) and 0.0147 AP$_S$, and both 512-image halves lost AP. Its
median application time was 1.1--1.4\,ms lower across three paired
runs (41.3--41.9 versus 42.4--43.2\,ms), but its 95th-percentile time
was not. Small-object recall at IoU 0.5 fell by 1.55 points (paired
image bootstrap 95\% interval $[-2.98, -0.19]$). The router was
therefore rejected. In the same units, the scheduling steps of
\cref{tab:ladder} save 15--28\,ms of median per-image latency,
depending on the starting pipeline. The router's ${\approx}1.2$\,ms is
at most one eleventh of that. Latency rather than period is compared because
the router was evaluated in a sequential pipeline, so its saving does
not compose with an overlapped schedule.

\section{QNN userspace notes}
\label{app:qnn}

Two findings about the bundled \texttt{qai-appbuilder}/QNN userspace
(QAIRT 2.50.40 wheel on QCM6490, V68 HTP) shaped the HTP campaign.
First, it is not safe under concurrent inference. Calls from
two contexts at once wedge the DSP session after a few hundred calls,
leaving an empty output or a session that never returns; retries and
context rebuilds do not recover it. This applies only to the bundled
userspace as shipped. The vendor SDK advertises multi-context support,
which was not exercised. The pipeline therefore holds exactly one
context and issues every NPU call from a single thread. Overlap comes
from pipelining decode, input fill and host NMS \emph{around} the
serialized NPU calls. Under that discipline the campaign ran
without a failed call.

Second, an earlier version of the evaluation script contained a harness
bug. It passed the input buffer to the appbuilder wrapper unflattened,
and the wrapper's reshape helper caught the resulting broadcast error
and printed it on every call, inside the timed region. This cost
${\approx}1$\,ms per call at the shipped governor and more at pinned
low clocks. Predictions were unaffected, because the native call
received the correct buffer. After the call was fixed, the entire HTP
campaign was rerun, and only the corrected numbers are reported. The
fix moved the shipped-governor rungs by ${\approx}1$\,ms per call, the
\texttt{performance} ladder by up to 13\% (the pinned core now spends
its cycles on dispatch instead of printing), and the 806\,MHz probe
point by 8\,ms per call. It therefore mattered most where the paper
examines clock sensitivity.
